\begin{frame}{Metodologia de validação}
  \small
  \begin{columns}[T,onlytextwidth]
    \begin{column}{0.31\textwidth}
      \begin{block}{1. Análise teórica}
        \begin{itemize}
          \item equações nodais e modelo ideal do ampop;
          \item resolução simbólica com SymPy;
          \item Bode entre $0{,}01$ e $10^5\,\mathrm{Hz}$;
          \item transformada inversa de Laplace.
        \end{itemize}
      \end{block}
    \end{column}
    \begin{column}{0.31\textwidth}
      \begin{block}{2. Simulação}
        \begin{itemize}
          \item duas configurações no LTspice;
          \item alimentação de $\pm10\,\mathrm{V}$;
          \item análise AC com entrada unitária;
          \item excitações transitórias de $0{,}8$ e $5\,\mathrm{V}$.
        \end{itemize}
      \end{block}
    \end{column}
    \begin{column}{0.31\textwidth}
      \begin{block}{3. Ensaio em bancada}
        \begin{itemize}
          \item conferência dos componentes;
          \item montagem em protoboard;
          \item gerador e osciloscópio de dois canais;
          \item medição de amplitudes, corte e tempos.
        \end{itemize}
      \end{block}
    \end{column}
  \end{columns}

  \begin{alertblock}{Grandezas de comparação}
    \centering
    $|H(j\omega)|=V_o/V_i$
    \qquad
    $V_o(f_c)=V_o(0)/\sqrt{2}$
    \qquad
    $t_r=t_{90}-t_{10}$.
  \end{alertblock}
\end{frame}
